\documentclass[11pt,letterpaper]{article}
\usepackage[T1]{fontenc}
\usepackage{lmodern,amsmath,amssymb,amsthm,microtype}
\usepackage[letterpaper,margin=1in]{geometry}
\usepackage[hidelinks]{hyperref}
\setlength{\emergencystretch}{2em}
\newtheorem{theorem}{Theorem}[section]
\newtheorem{lemma}[theorem]{Lemma}
\newtheorem{proposition}[theorem]{Proposition}
\newtheorem{corollary}[theorem]{Corollary}
\theoremstyle{remark}\newtheorem{remark}[theorem]{Remark}
\newcommand{\F}{\mathcal F}
\newcommand{\nc}[2]{#1\mathrm N#2}
\newcommand{\sw}{\operatorname{swap}}
\newcommand{\del}{\Sigma}
\newcommand{\maj}{\operatorname{maj}}
\renewcommand{\thesection}{S\arabic{section}}
\title{Certificate Supplement to\\
Improved Upper Bounds on the Size of Condorcet Domains}
\author{Pavlos Moros \and S\o ren Riis}
\date{Local review draft, 29 October 2026}
\hypersetup{pdftitle={Certificate Supplement: Improved Upper Bounds on the Size of Condorcet Domains},pdfauthor={Pavlos Moros and Soren Riis}}
\begin{document}
\maketitle
\section{Scope and notation}
This supplement records the finite lemma, its certificate and exact
reproduction instructions. The general proof remains in the main paper.
Certificate files are unchanged from the frozen 29 October source;
no finite generation or audit was rerun for this editorial preparation.

Let $\F_k$ consist of the nonempty families of orders on $k$ alternatives
defined by one extreme never-condition on each triple and arbitrary
fixed pair comparisons. A condition $\nc xp$ means that $x$ never has
position $p$ on its triple; extreme means $p=1$ or $p=3$.
Write $A(H)$ for the number of enabled adjacent swaps, counted over
orders and positions, so each undirected edge contributes twice.
The definitions below reproduce the exact finite argument used in the
manuscript, with additional enumeration and selector details.

\section{A finite window inequality}\label{sec:finite}

For $H\in\F_k$, its \emph{bottom slices} group orders by their last
alternative and delete it; top slices group by the first alternative.
Every nonempty slice belongs to $\F_{k-1}$. Write $a_j(H)$ for the
number of orders whose swap at position $j$ is enabled. On seven labels put
\begin{equation}\label{eq:cost}
 C(H)=\frac{a_1+2a_2+2a_3+2a_4+2a_5+a_6}{10}.
\end{equation}
The weights sum to one. They are chosen so that each interior swap is
counted with total weight one when a seven-position window slides
through an order.

\begin{lemma}[Exact finite estimate]\label{lem:finite}
There is a rational function $q$ on the isomorphism types of $\F_6$ such
that, for every $J\in\F_7$,
\begin{equation}\label{eq:finite}
 C(J)\le\frac58|J|+
 \sum_{H\text{ bottom slice}}q(H)-\sum_{H\text{ top slice}}q(H).
\end{equation}
The rational table and exact verifiers are in
\texttt{ancillary/finite\_certificate/}. The finite statement includes
all fixed-pair constraints, with no common-axis assumption.
\end{lemma}
\begin{proof}
Here is the construction and the coverage argument for the finite check.
The certificate assigns per-order weights $\phi$ to five-label types.
For a six-domain $H$ let $b(H)$ and $t(H)$ be the sums of
$|L|\phi(L)$ over its bottom and top five-slices. Put
\[
 s(H)=\frac{12}{19}|H|+b(H)-t(H)-\frac15 A(H),\qquad \eta=\frac1{152}.
\]
There are specified exceptional sets $\mathcal T,\mathcal S$ of five and
nine six-label types. Put $p(H)=2\eta(r_t(H)-r_b(H))$ outside $\mathcal S$,
and $p(H)=\psi(H)$ on $\mathcal S$, with the rational values $\psi$ supplied
by the certificate. Set $\theta_b=s/2+p$ and $\theta_t=s/2-p$.
The finite facts are as follows. On $\mathcal T$
both split values are zero. Outside $\mathcal T\cup\mathcal S$,
\begin{equation}\label{eq:receivers}
 \theta_b(H)\ge\frac\eta2|H|+4\eta r_t(H),\qquad
 \theta_t(H)\ge\frac\eta2|H|+4\eta r_b(H).
\end{equation}
Here $r_b,r_t$ count specified five-slices described below. A further
finite check proves
\begin{equation}\label{eq:lambda}
 \Lambda(J):=\sum_{\text{bottom }H}\theta_b(H)
             +\sum_{\text{top }H}\theta_t(H)\ge\eta|J|
\end{equation}
whenever $J$ has a slice in $(\mathcal T\setminus\{R_8\})\cup\mathcal S$.
All types, weights and slice selectors are explicitly defined in the
ancillary data; they are not fitted during verification.

For completeness, the one remaining exceptional type is
\[
 R_8=\{524130,524310,534120,534210,
        542130,542310,543120,543210\}.
\]
It has a common first and last alternative. In the regenerated five-label
class map the two selector types below have indices 866 and 41, respectively. The quantity $r_b(H)$ counts
bottom five-slices of the type $R_8$ with its first alternative deleted;
$r_t(H)$ counts top five-slices of the type with its last deleted.
Both counts vanish on $R_8$.
If a bottom $R_8$ slice occurs, its eight orders share their first
alternative, so they enter one top receiver slice and contribute one
to its $r_b$. That receiver is not $R_8$. Distinct bottom slices give
distinct receiver/label pairs. Each contributes $4\eta$, exactly paying
the missing baseline $8\eta/2$. Top $R_8$ slices are symmetric.
In the remaining cases all non-$R_8$ slices lie outside
$\mathcal T\cup\mathcal S$, so~\eqref{eq:receivers} applies to every receiver.
It proves~\eqref{eq:lambda} in every case not handled by the exceptional finite check.

Define
\[
 q(H)=\frac{b(H)+t(H)}2+p(H).
\]
In a seven-domain, a top slice of a bottom slice is the same conditional
five-domain as the corresponding bottom slice of a top slice.
These middle terms cancel. Also the sum of $A(H)$ over all six-slices
is $10C(J)$. Consequently
\[
 \Lambda(J)=\frac{12}{19}|J|-C(J)
       +\sum_{\text{bottom }H}q(H)-\sum_{\text{top }H}q(H).
\]
Now~\eqref{eq:lambda} and $12/19-1/152=5/8$ give~\eqref{eq:finite}.

We finally explain why the finite checks are exhaustive.
All $128\,400$ labelled five-domains form $1\,160$ relabelling types.
Extend their defining laws over the ten new triples and five new pairs.
This gives $14\,105\,084$ nonempty six-extensions of five-type
representatives and $94\,830$ feature rows. The extension count is not a
count of distinct labelled six-domains.
Every six-domain occurs up to relabelling: restrict its \emph{defining
conditions} to five labels, rather than taking its projection.
Keeping every domain on the thirteen exceptional feature rows gives precisely the
fourteen required types. Two of the five types in $\mathcal T$ share one feature row.
Their representatives and reversals, omitting
$R_8$, give eighteen seven-label seed families. Requiring a prescribed
endpoint six-slice and intersecting all compatible law and pair options
produces $14\,248$ seed--domain cases (before deduplication across seeds)
and $18\,316$ feature rows, covering both
orientations. Every required inequality is verified by clearing the rational
denominators and checking a nonnegative integer residual. A separate
tuple-set checker recomputes every saved row from its witness orders.
The written coverage and receiver arguments are given above. The separate
Lean project in the repository also formalises these steps and the passage
to arbitrary $n$; its verification report records the checked scope.
\end{proof}


\section{Reproducing the finite checks}
The repository directory \texttt{ancillary/finite\_certificate/} contains
standard Python~3
scripts, two C++17 generator sources, rational certificate data, saved
feature rows and tuple witnesses. No optimizer, network access or
nonstandard Python package is needed. Run commands from that directory.

\paragraph{Saved-data replay.}
\begin{verbatim}
sh run_joint_checks.sh
\end{verbatim}
This checks the rational residuals on saved rows, recomputes features
from all $94\,830+18\,316=113\,146$ tuple witnesses, checks the
compressed potential and scalar transfers, and runs supplementary small
charging-map cases. Successful completion has exit status zero. It
checks the supplied data, rather than reconstructing the family coverage.

\paragraph{Complete regeneration.}
With a C++17 compiler named \texttt{c++} available, use a new output path:
\begin{verbatim}
python3 run_fresh_regeneration.py /path/to/new/output
\end{verbatim}
The driver refuses to overwrite an existing output directory. It starts
from generator sources, coefficient data and exceptional keys, not saved
feature rows or witnesses. It regenerates the five-label families,
six-extensions and required seven-label families, checks completion flags,
runs the exact and tuple verifiers, and compares regenerated feature
tables with the supplied tables. Its success receipt is
\texttt{FULL\_REGENERATION\_RECEIPT.json}, with result
\texttt{COMPLETE FRESH REGENERATION PASS}. Each command and return code
is recorded there. Witness order can depend on hash iteration; feature
rows are compared exactly, and every regenerated witness is checked.

\paragraph{Expected counts.}
The units are different and should not be interchanged:
\begin{itemize}
\item $128\,400$ labelled five-families and $1\,160$ relabelling types;
\item $14\,105\,084$ six-extensions and $94\,830$ six-feature rows;
\item $397$ selected sets, giving fourteen exceptional types on thirteen
feature rows;
\item eighteen seven-label seeds, $14\,248$ seed--domain cases before
cross-seed deduplication, and $18\,316$ seven-feature rows.
\end{itemize}
Both orientations are covered. The generic and $R_8$ receiver cases are
the hand arguments in Section~\ref{sec:finite}; they are not replaced by
a success flag. The separate Lean project formalises the coverage proof and the passage
from seven positions to arbitrary $n$; it is documented in the repository.

\section{Verification}
The repository records the prior complete finite regeneration and a fresh
saved-data replay. Lean sources, documentation and verification scope are
provided separately through the GitHub repository and archived release
referenced in the main paper.
\end{document}
